\ssscp{\tsc{Peripheral Babar}}{04-06-05-01}
There is evidence for a group containing West Damer,
Dawera-Daweloor, Dai, and North Babar.
We call this group \tsc{Peripheral Babar} 
to avoid confusion with the North Babar language [bcd]
which is a member of this group.
This group was first suggested by \citet{vanEngelenhoven2010}
in his discussion of the relationships of West Damer.

There is one change which applies to all of these languages
and four partially overlapping sound changes which group these
languages together in a wave-like pattern.
These sound changes are given in \qf{04-26} below,
with a wave diagram showing the internal relations of the \tsc{Peripheral
Babar} languages given in \frf{04-05}. Each line in this diagram
represents a different sound change.
These changes are discussed further below.

{\wg\st{6pt}\tblr{>{\hspace{-\tabcolsep}}l<{\hspace{-\tabcolsep}}>{\hspace{-\tabcolsep}\enspace}llllll}
\tbx{04-26}	&	\mc{6}{l}{\tsc{Peripheral Babar} sound changes}	\\
						&	a.	&	*s > **ʈ				&	>	&	\it{d}				&							&all \tsc{Peripheral Babar} languages	\\
						&	b.	&	*V\tsc{[+high]}	&	>	&	C							&	/{\gap}-V		&	West Damer, Dawera-Daweloor, Dai	\\
						&	c.	&	*ŋ							&	>	&	\it{ʔ >} {\0}	&							&Dawera-Daweloor, Dai, North Babar	\\
						&	d.	&	*l							&	>	&	\it{n}				&							&Dai, North Babar	\\
						&	e.	&	*V							&	>	&	{\0}					&	/{\gap}{\#}	&	Dawera-Daweloor, Dai	\\
\end{tabular}\mbox{}\\}


\begin{figure}\centering
	\includegraphics[width=0.75\columnwidth]{Babar-wave.pdf}
	\fcp{\tsc{Peripheral Babar} wave diagram}{04-05}
\end{figure}

The first change is *s > **ʈ > \it{d},
which was discussed above.
This is the only change which all \tsc{Peripheral Babar} languages share in common.
The second change is fortition of high vowels in root-final position
before a vocalic suffix.
This change affects West Damer, Dawera, and Dai.
West Damer and Dai have *i > \it{s},
while Dawera has *i > \it{sy} where /sy/ → [s] {\tl} [si] {\tl} [ʃ] {\tl} [ʃi].
The usual reflexes of *u are: *u > \it{k} in Dawera,
*u > \it{ʔ} in Dai,
and *u > \it{h} in West Damer.
West Damer /h/ has a voiceless
dorsal fricative allophone [x] in \citet{Chlenova2008}.\footnote{
		\citet[143]{Chlenova2008}
		has ‹\it{ch}› as a variant of ‹\it{h}›.
		The data in \citet{Chlenova2008} was handwritten by their consultants
		in the 1960s on the model of the former Indonesian orthography
		in which ‹\it{ch}› represents a voiceless velar fricative [x].}
Examples of this vowel fortition are given in \trf{04-49}.

\begin{table}\tcp{\tsc{Peripheral Babar} high vowel fortition /{\gap}-V}{04-49}
\begin{adjustbox}{width=\textwidth}
	\begin{threeparttable}\st{2pt}
	\begin{tabular}{llllllll}\lsptoprule
local gl.	&	{\hr}stone	&	{\hr}hair	&	{\hr}dog	&	{\hr}louse	&	{\hr}breast	&	{\hr}bone	&	{\hr}fruit bat\\
PMP	&	\hr*bat\tbr{u}	&	\hr*bul\tbr{u}	&	\hr*as\tbr{u}	&	\hr*kut\tbr{u}	&	\hr*sus\tbr{u}	&	\hr*duR\tbr{i}	&	\hr*panik\tbr{i}\\
P.-Babar	&	**wat\tbr{u}	&	**wul\tbr{u}	&	**aʈ\tbr{u}	&	**ut\tbr{u}	&	**dud\tbr{u}	&	**ru\tbr{i}	&	**nik\tbr{i}\\ \midrule
W. Damer	&	{\hr\hr}wot\tbr{h}-o	&	{\hr\hr}wul{\tl}wul\tbr{h}-o	&	{\hr\hr}ot\tbr{h}-o	&	{\hr\hr}ut\tbr{h}-o	&	{\hr\hr}dut\tbr{h}-o	&	{\hr\hr}ru\tbr{s}-o	&	{\hr\hr}ni\tbr{s}-o\\
Dawera	&	{\hr\hr}wat\tbr{k}-ol	&	{\hr\hr}wul\tbr{k}-ol	&	{\hr\hr}ad\tbr{k}-ol	&	{\hr\hr}ot\tbr{k}-ol	&	{\hr\hr}dud\tbr{k}-ol	&	{\hr\hr}ru\tbr{sy}-ol	&	{\hr\hr}lik\tbr{sy}-ol\\
Dai	&	{\hr\hr}wat\tbr{ʔ}-on	&	{\hr\hr}wun\tbr{ʔ}-on	&	{\hr\hr}ad\tbr{h}-on	&	{\hr\hr}ot\tbr{ʔ}-on	&	{\hr\hr}dud\tbr{ʔ}-on	&	{\hr\hr}ru\tbr{s}-on	&	{\hr\hr}nit\tbr{s}-on\su{‡}\\
N. Babar	&	{\hr\hr}waʔa\tbr{y}-ai	&	{\hr\hr}wun\tbr{y}-ai	&	{\hr\hr}aduw\tbr{i}-ai	&	{\hr\hr}ot\tbr{i}-ai	&	{\hr\hr}dut-niɛ\su{†}	&	{\hr\hr}ru-n-ai	&	{\hr\hr}niʔ\tbr{i}-ai\\
		\lspbottomrule
	\end{tabular}
		\begin{tablenotes}\tnsize
			\item[†]	{North Babar \it{dutniɛ} ‘breast’ has high vowel metathesis
								across morpheme boundaries; \it{duti + -nɛ → dutniɛ}.}
			\item[‡]	{Dai \it{nits-on} ‘fruit bat’ has *k > \it{t} before an alveolar
								consonant, also seen in *hikan > \it{itn-on} ‘fish’.}
		\end{tablenotes}
	\end{threeparttable}
\end{adjustbox}
\end{table}

Additional examples of *i > \it{s} in West Damer include *tal\tbr{i}h
> \it{tol\tbr{s}-o} ‘rope’, and *hapuy > **a\tbr{i} > \it{o\tbr{s}-o} ‘fire’.
Additional examples of *i > \it{s} in Dawera
and Dai in words inherited from PMP are scarce,
but there are regionally distributed cognate sets which establish
the regularity of the pattern.
Such examples include **pe\tbr{i} > Dawera \it{pe\tbr{sy}} [peʃ],
Dai \it{pei\tbr{s}-on} ‘coconut’, **a\tbr{i} > Dawera \it{a\tbr{sy}-ol}
[ˈaʃiol], Dai \it{ai\tbr{s}-on} ‘boat’, and **e\tbr{i} > Dawera \it{e\tbr{sy}-ol}
[ˈeʃiol], Dai \it{ɛ\tbr{s}-on} ‘person’.
Each of these can be compared with North Babar \it{pei-ai} [ˈpeiya] ‘coconut’,
\it{ai-ai} [ˈayai] ‘boat’, and \it{ei-ai} [ˈejai] ‘person’, which retain earlier final **i.

Vowel fortition probably originally arose through fortition of
phonetic glides.
The way this probably happened is illustrated in \qf{04-27} below for
Dawera *paniki > \it{niksy-ol} ‘fruit bat’
and *batu > \it{watk-ol} ‘stone’.

{\tblr{>{\hspace{-\tabcolsep}}l<{\hspace{-\tabcolsep}}>{\hspace{-\tabcolsep}\enspace}l@{ }l@{ }l@{ }l@{ }l@{ }l@{ }l@{ }l@{ }l@{ }l@{ }l}
\tbx{04-27}	&	*paniki	&	>	&	**niki-ol	&	>	&	**nikiy-ol	&	>	&	**nikiʤ-ol	&	>	&	**nikisy-ol	&	>	&	\it{niksy-ol}	\\
						&	*batu		&	>	&	**watu-ol	&	>	&	**watuw-ol	&	>	&	**watug-ol	&	>	&	**watuk-ol	&	>	&	\it{watk-ol}	\\
\end{tabular}\mbox{}\\}

Evidence for this process comes from \citet{Taber1993} who records
reduced vowels in some words when no suffix occurs.
Two examples are *kutu > Dawera \it{otk} [ˈotə̆kʰ] ‘louse’,
and *batu > \it{watk} [ˈwatək] ‘stone’.
The vowels in these words are probably retentions of the original
final vowel represented at the penultimate stage in \qf{04-27}.\footnote{
		An anonymous reviewer suggests that these vowels could be due to epenthesis.
		While this is possible,
		simplicity favours treating them as retentions.
		Under the epenthesis analysis we need another step in \qf{04-27}.} 
Another piece of evidence for this process is that some words
have variants with and without the final consonant in different phrases.
This could be the result of these words originally occurring
with or without suffixes in different contexts.
One example is *babuy > **wawi > **wowi > West Damer \it{wows-o}
‘pig’, \it{wowi kmirsoni} ‘the fat pig’,
and \it{wow daroni} ‘wild pig’ \citep[175]{Chlenova2008}.
Although the best analysis of the limited examples is unclear,
this data bears a striking resemblance to the more well described
word-shape alternations in other languages of the region such
as Leti \citep{vanEngelenhoven2004} and Luang \citep{TaberTaber2015}.

While vowel fortition may have arisen only before suffixes,
it appears to have become lexicalised to a certain degree.
While there are examples of alternations,
there are also many examples in which the final consonant occurs
without any following vowel initial suffix.
One example from West Damer is *asu > \it{oth purani} ‘wild dog’.
Examples from Dawera include *batu > \it{watk-re} ‘stone-\tsc{pl}’,
and \it{watk pus{\tl}pus} ‘small stones’ \citep[148]{Chlenova2002},
as well as the isolated Dawera words *kutu > \it{otk} [ˈotə̆kʰ]
‘louse’, and *batu > \it{watk} [ˈwatək] ‘stone’ in \citet{Taber1993}.

The changes affecting PMP *ŋ affect all \tsc{Peripheral Babar}
languages except West Damer.
Dawera and North Babar have *ŋ > {\0} in all positions,
while Dai has initial *ŋ > {\0}, but medial *ŋ > \it{ʔ}.
Medial *ŋ > \it{ʔ} in Dai is probably an intermediate stage before loss.
Examples of *ŋ > **ʔ > {\0} in \tsc{Peripheral Babar} are given in \trf{04-50}.
West Damer has regionally common *ŋ > \it{n}.

\begin{table}\tcp{\tsc{Peripheral Babar} *ŋ > **ʔ > Ø}{04-50}
\begin{adjustbox}{width=\textwidth}
	\begin{threeparttable}\st{2pt}
	\begin{tabular}{llllll@{}ll}\lsptoprule
local gl.	&	{\hr}name	&	{\hr}tooth	&	{\hr}ear	&	{\hr}wind	&	{\hr}swim	&	{\hr}branch	&	{\hr}garden\\
PMP	&	\hr*\tbr{ŋ}ajan	&	\hr*\tbr{ŋ}isi	&	\hr*tali\tbr{ŋ}a	&	\hr*haŋin	&	\hr*naŋuy	&	\hr*saŋa	&	\\
P.--Babar	&	**\tbr{ŋ}an	&	**\tbr{ŋ}iʈ	&	**tli\tbr{ŋ}a	&	**a\tbr{ŋ}in	&	**wena\tbr{ŋ}ui	&	**ʈa\tbr{ŋ}a	&	**ti\tbr{ŋ}a\su{Δ}\\	\midrule
W. Damer	&	{\hr\hr}\tbr{n}on-o	&	\hr\cg{(nin-o)}	&	{\hr\hr}tli\tbr{n}a	&	{\hr\hr}o\tbr{n}yn-o	&	{\hr\hr}ye-nu{\tl}nu\tbr{n}i	&	{\hr\hr}do\tbr{n}a	&	{\hr\hr}ti\tbr{n}a\\
Dawera	&	\hr\cg{(wuletal-ol)}	&	{\hr\hr}id-ol	&	{\hr\hr}kli-ol	&	{\hr\hr}aili-ol	&	{\hr\hr}welaiy, welauy	&	{\hr\hr}da-ol	&	{\hr\hr}ti-ol\\
Dai	&	\hr\cg{(wnyɛtan)}	&	{\hr\hr}id-on	&	{\hr\hr}nni\tbr{ʔ}-on	&	\hr\cg{(suʔ-on)}	&	{\hr\hr}n-wɛna\tbr{ʔ}os\su{‡}	&		&	\\
N. Babar	&	{\hr\hr}an-ai	&	{\hr\hr}it-ne	&	{\hr\hr}kli-ne\su{†}	&	{\hr\hr}ani-ai\su{†}	&	{\hr\hr}n-wɛnai	&		&	\\
		\lspbottomrule
	\end{tabular}
		\begin{tablenotes}\tnsize
			\item[†]	{\citet{Taber1993} transcribes ‹nn› for North Babar [ˈklinnɛ]
								‘ear’ and [ˈanniai] ‘wind’.
								This could attest *ŋ > \it{n},
								but there are also indications that North Babar has optional
								phonetic lengthening of /n/.
								Examples of [nn] in words for which only one \it{n} is expected
								on etymological grounds include: *mata > *mata-ne > [ˈmatənnɛ]
								‘eye-\tsc{3sg}’, and *ina > [ˈɛnne] ‘mother’.}
			\item[‡]	{Dai \it{n-wenaʔos} ‘swim’ has final *i > \it{s} from intermediate **wenaʔoi.}
			\item[Δ]	{Proto-Babar **tiŋa ‘garden’ has widespread regional cognates
								including: Kisar \it{kirna}, Roma \it{tiran} ‘garden’, Tetun \it{tinan} ‘year’,
								Termanu \it{tina} ‘dry field which is replanted every year’,
								and Welaun (\tsc{Central Timor}) \it{tikan} ‘year’.
								These point to earlier **tiŋaR,
								as suggested by \citet[285]{Mills2010}.}
		\end{tablenotes}
	\end{threeparttable}
\end{adjustbox}
\end{table}

There are two other sound changes found in more than one \tsc{Peripheral
Babar} language: Dai and North Babar share *l > \it{n},
while Dai and Dawera have lost most final vowels which did not undergo fortition.
Loss of final vowels also occurs to a lesser extent in West Damer and North Babar.
Examples of *l > \it{n} in Dai
and North Babar, as well as *V > {\0} /{\gap}{\#} in Dawera-Daweloor
and Dai are given in \trf{04-51}.

\begin{table}\tcp{\tsc{Peripheral Babar} *l and *V /C{\gap\#}}{04-51}
\begin{adjustbox}{width=\textwidth}
	\st{2pt}\begin{tabular}{llll|llll}\lsptoprule
local gl.	&	{\hr}head	&	{\hr}sun	&	{\hr}five	&	{\hr}vomit	&	{\hr}eye	&	{\hr}\tsc{1pl.incl}	&	{\hr}house\\
PMP	&	\hr*qu\tbr{l}u	&	\hr*qa\tbr{l}əjaw	&	\hr*\tbr{l}ima	&	\hr*um-ut\tbr{a}q	&	\hr*mat\tbr{a}	&	\hr*kit\tbr{a}	&	\hr*Rum\tbr{a}q\\
P.-Babar	&	**u\tbr{l}u	&	**\tbr{l}aya	&	**\tbr{l}ima	&	**mut\tbr{a}	&	**mat\tbr{a}	&	**kit\tbr{a}	&	**um\tbr{a}\\ \midrule
W. Damer	&	{\hr\hr}ulk-ona	&	{\hr\hr}law-oni	&	{\hr\hr}wilimo	&	{\hr\hr}wamuto	&	{\hr\hr}mota	&	{\hr\hr}itito	&	{\hr\hr}uma\\
Dawera	&	\hr\cg{(radl-ol)}	&	{\hr\hr}ley-ol	&	{\hr\hr}lim	&	{\hr\hr}lɛ-mut	&	{\hr\hr}mat-ol	&	{\hr\hr}it	&	{\hr\hr}um-ol\\
Dai	&	{\hr\hr}u\tbr{n}ʔ-on	&	{\hr\hr}\tbr{n}e-on	&	{\hr\hr}\tbr{n}im	&	{\hr\hr}n-mut	&	{\hr\hr}mat-on	&	{\hr\hr}it	&	{\hr\hr}um-on\\
N. Babar	&	{\hr\hr}u\tbr{n}-ni	&	{\hr\hr}\tbr{n}eya	&	{\hr\hr}\tbr{n}imo	&	{\hr\hr}in-moʔo	&	{\hr\hr}matə-nɛ	&	{\hr\hr}ita	&	{\hr\hr}om-ai\\
		\lspbottomrule
	\end{tabular}
\end{adjustbox}
\end{table}

Some other changes shared among \tsc{Peripheral Babar} also occur
in \tsc{South Babar} languages.
Dawera has *n > \it{l},
and North Babar has *t > \it{k} in certain environments.
Examples of *n > \it{l} in Dawera are given in \trf{04-52}.

\begin{table}\tcp{\tsc{Peripheral Babar} *n}{04-52}
\begin{adjustbox}{width=\textwidth}
	\st{2pt}\begin{tabular}{ll@{}lll@{}lll}\lsptoprule
local gl.	&	{\hr}child	&	{\hr}fish	&	{\hr}bird	&	{\hr}six	&	{\hr}star	&	{\hr}man	&	{\hr}sand\\
PMP	&	\hr*a\tbr{n}ak	&	\hr*hika\tbr{n}	&	\hr*ma\tbr{n}uk	&	\hr*ə\tbr{n}əm	&	\hr*bituqə\tbr{n}	&	\hr*maRuqa\tbr{n}ay	&	\hr*qə\tbr{n}ay\\
P.-Babar	&	**a\tbr{n}	&	**ika\tbr{n}	&	**ma\tbr{n}u	&	**\tbr{n}emo	&	**tu\tbr{n}	&	**ma\tbr{n}i	&	**eni\\ \midrule
W. Damer	&	\hr\cg{(kəmemey-o)}	&	{\hr\hr}yeny-o	&	{\hr\hr}munw-o	&	{\hr\hr}wyɛ-namo	&	{\hr\hr}ton-o	&	\hr\cg{(ulh-o)}	&	{\hr\hr}ɛn-o\\
Dawera	&	{\hr\hr}a\tbr{l}-ol	&	{\hr\hr}i\tbr{l}	&	{\hr\hr}ma\tbr{l}-ol	&	{\hr\hr}\tbr{l}ɛm	&	{\hr\hr}to\tbr{l}-ol	&	{\hr\hr}mu\tbr{l}i-ol	&	{\hr\hr}ɛ\tbr{l}i-ol\\
Dai	&	{\hr\hr}an-on	&	{\hr\hr}itn-on	&	{\hr\hr}manʔ-on	&	{\hr\hr}nɛm	&	\hr\cg{(tnorn-on)}	&	\hr\cg{(es-on)}	&	{\hr\hr}ɛny-on\\
N. Babar	&	{\hr\hr}an-ai	&	{\hr\hr}in-ai	&	{\hr\hr}mani-a	&	{\hr\hr}nemo	&	{\hr\hr}toʔoni-a	&	{\hr\hr}mɛny-ai	&	{\hr\hr}ɛni-a\\
		\lspbottomrule
	\end{tabular}
\end{adjustbox}
\end{table}

\begin{table}\tcp{\tsc{Peripheral Babar} *t}{04-53}
\begin{adjustbox}{width=\textwidth}
	\st{2pt}\begin{tabular}{ll@{}l@{}llll@{}l@{}l}\lsptoprule
local gl.	&	{\hr}\tsc{1pl.excl}	&	{\hr}deaf	&	{\hr}burn	&	{\hr}knee	&	{\hr}egg	&	{\hr}soil	&	{\hr}three	&	{\hr}black\\
PMP	&	\hr*ki\tbr{t}a	&	\hr*\tbr{t}ilu	&	\hr*\tbr{t}unu	&	\hr*\tbr{t}uhud	&	\hr*qa\tbr{t}əluR	&	\hr*\tbr{t}anəq	&	\hr*\tbr{t}əlu	&	\hr*ma-qi\tbr{t}əm\\
P.-Babar	&	**i\tbr{t}a	&	**\tbr{t}ilu	&	**\tbr{t}uni	&	**\tbr{t}ur	&	**\tbr{t}eli	&	**\tbr{t}ani	&	**\tbr{t}eli	&	**me\tbr{t}em\\ \midrule
W. Damer	&	{\hr\hr}iti\tbr{t}o	&	{\hr\hr}k\tbr{t}ili	&	{\hr\hr}k\tbr{t}uni	&	{\hr\hr}\tbr{t}ur-o	&	{\hr\hr}diw-\tbr{t}alh-o	&	\hr\cg{(tanah)}	&	{\hr\hr}wyɛ-\tbr{t}eli	&	{\hr\hr}mɛt{\tl}mɛ\tbr{t}m-o\\
Dawera	&	{\hr\hr}i\tbr{t}	&	{\hr\hr}\tbr{t}ilk-ol	&	{\hr\hr}\tbr{t}ut	&	{\hr\hr}\tbr{t}ur-ol	&	{\hr\hr}man \tbr{t}ilk	&	\hr\cg{(waw)}	&	{\hr\hr}\tbr{t}ɛl	&	{\hr\hr}mɛ\tbr{t}me-l\\
Dai	&	{\hr\hr}i\tbr{t}	&	{\hr\hr}\tbr{t}inʔ-on	&	{\hr\hr}\tbr{t}unɛn	&	{\hr\hr}\tbr{t}ur-on	&	{\hr\hr}\tbr{t}enʔɛn-on	&	{\hr\hr}\tbr{t}any-on	&	{\hr\hr}tɛn	&	{\hr\hr}n-mɛtm\\
N. Babar	&	{\hr\hr}i\tbr{t}a	&	{\hr\hr}n-\tbr{t}ini	&	{\hr\hr}\tbr{t}uni	&	{\hr\hr}\tbb{k}or-ne	&	{\hr\hr}diw \tbb{k}eny-ai	&	{\hr\hr}\tbb{k}any-ai	&	{\hr\hr}\tbb{k}ini	&	{\hr\hr}n-mɛ\tbb{ʔ}em-o\\
		\lspbottomrule
	\end{tabular}
\end{adjustbox}
\end{table}

North Babar has a split of *t > \it{t, k, ʔ}.
\citet{Zobel2018} tentatively posits that *t = \it{t} adjacent
to a high vowel *i or *u,
but *t > \it{k, ʔ} elsewhere.
Under his hypothesis, *u > \it{o} preceded *t > \it{k, ʔ}.
This hypothesis works for much of the data, but exceptions occur.
Examples in the environments proposed by \citet{Zobel2018}
are given in \trf{04-53}.
Examples which occur in other environments include: *ma\tbr{t}ay
> \it{n-ma\tbr{ʔ}i} ‘die’, *\tbr{t}aqi > \it{\tbr{t}ei-ei} ‘faeces’,
*ma\tbr{t}a > \it{ma\tbr{t}ə-nɛ} ‘eye’,
and *hu\tbr{t}ək > \it{o\tbr{t}ɛ-nə nani} ‘brain’.
